\documentclass{article}
\usepackage{amsmath,amssymb,amsthm}
\usepackage{algorithm}
\usepackage{algpseudocode}
\usepackage{enumitem}
\usepackage{hyperref}
\usepackage{bm}
\usepackage{geometry}
\geometry{margin=1in}
\newtheorem{theorem}{Theorem}
\newtheorem{lemma}{Lemma}
\newtheorem{assumption}{Assumption}
\newcommand{\prox}{\operatorname{prox}}
\newcommand{\grad}{\nabla}
\newcommand{\dist}{\operatorname{dist}}
\newcommand{\R}{\mathbb{R}}
\newcommand{\E}{\mathbb{E}}

\begin{document}

\section*{Algorithm Name}
\textbf{Proximal Gradient Method (PGM) for a smooth non‑convex loss $f$ and a (convex) regulariser $g$.}

\section*{Mathematical Formulation}
Let $f:\R^{d}\to\R$ be differentiable with $L$‑Lipschitz continuous gradient and let $g:\R^{d}\to\R\cup\{+\infty\}$ be proper, closed and convex (so that $\prox_{\eta g}$ is well defined).  
Define the composite objective
\[
\Phi(x)\;:=\;f(x)+g(x),\qquad x\in\R^{d}.
\]
For a stepsize $\eta>0$ the iteration reads
\begin{equation}
\boxed{\;
x_{k+1}\;=\;\prox_{\eta g}\!\bigl(x_{k}-\eta\,\grad f(x_{k})\bigr),\qquad k=0,1,2,\dots
\;}
\label{eq:update}
\end{equation}
The \emph{gradient mapping} (a surrogate for the stationary condition) is
\[
G_{\eta}(x)\;:=\;\frac{1}{\eta}\Bigl(x-\prox_{\eta g}\!\bigl(x-\eta\,\grad f(x)\bigr)\Bigr).
\tag{1}
\]

\section*{Pseudocode}
\begin{algorithm}[H]
\caption{Proximal Gradient Method}
\begin{algorithmic}[1]
\Require Initial point $x_{0}\in\R^{d}$, stepsize $\eta\in(0,1/L)$, max iter $K$.
\For{$k=0$ \textbf{to} $K-1$}
    \State $u_{k}\gets x_{k}-\eta\,\grad f(x_{k})$ \Comment{gradient step}
    \State $x_{k+1}\gets \prox_{\eta g}(u_{k})$ \Comment{proximal step}
\EndFor
\State \Return $\{x_{k}\}_{k=0}^{K}$.
\end{algorithmic}
\end{algorithm}

\section*{Dry Run Example}
Consider the scalar problem $f(w)=(w-3)^{2}$, $g\equiv0$ (hence $\prox_{\eta g}(z)=z$).  
Take $w_{0}=0$, stepsize $\eta=0.1$ and run three iterations.

\begin{enumerate}[label=\textbf{Iter \arabic*:},leftmargin=*]
\item \textbf{Gradient:} $\grad f(w_{0})=2(w_{0}-3)=-6$.
\newline
\textbf{Update:} $w_{1}=w_{0}-\eta\grad f(w_{0})=0-0.1(-6)=0.6$.
\newline
\textbf{Objective:} $\Phi(w_{1})=f(w_{1})=(0.6-3)^{2}=5.76$.
\item \textbf{Gradient:} $\grad f(w_{1})=2(0.6-3)=-4.8$.
\newline
\textbf{Update:} $w_{2}=0.6-0.1(-4.8)=1.08$.
\newline
\textbf{Objective:} $\Phi(w_{2})=(1.08-3)^{2}=3.6864$.
\item \textbf{Gradient:} $\grad f(w_{2})=2(1.08-3)=-3.84$.
\newline
\textbf{Update:} $w_{3}=1.08-0.1(-3.84)=1.464$.
\newline
\textbf{Objective:} $\Phi(w_{3})=(1.464-3)^{2}=2.365$.
\end{enumerate}
The objective decreases monotonically, illustrating the descent property proved below.

\newpage
\section*{Assumptions}
\begin{assumption}[Smoothness of $f$]\label{as:smooth}
$f$ is differentiable and its gradient is $L$‑Lipschitz:
\[
\|\grad f(x)-\grad f(y)\|\le L\|x-y\|,\qquad\forall x,y\in\R^{d}.
\]
\end{assumption}
\begin{assumption}[Convexity of $g$]\label{as:convex}
$g$ is proper, closed and convex, so that for any $\eta>0$ the proximal operator
\[
\prox_{\eta g}(z):=\arg\min_{x}\Bigl\{g(x)+\frac{1}{2\eta}\|x-z\|^{2}\Bigr\}
\]
is single‑valued and \emph{non‑expansive}:
\[
\|\prox_{\eta g}(u)-\prox_{\eta g}(v)\|\le\|u-v\|,\qquad\forall u,v\in\R^{d}.
\tag{2}
\]
\end{assumption}
\begin{assumption}[Stepsize]\label{as:stepsize}
The stepsize satisfies
\[
0<\eta<\frac{2}{L}.
\]
In the sequel we will use the more convenient bound $0<\eta<1/L$ which guarantees a strict descent.
\end{assumption}
\begin{assumption}[Existence of a lower bound]\label{as:bound}
The composite objective $\Phi$ is bounded below, i.e.
\[
\Phi^{\inf}:=\inf_{x\in\R^{d}}\Phi(x)>-\infty .
\]
\end{assumption}

\section*{Main Convergence Theorem}
\begin{theorem}[Convergence of the gradient mapping]\label{thm:convergence}
Let $\{x_{k}\}$ be generated by \eqref{eq:update} with a stepsize $\eta\in(0,1/L)$.  
Define the gradient mapping $G_{\eta}$ as in (1). Then for every $K\ge1$
\[
\min_{0\le k\le K-1}\|G_{\eta}(x_{k})\|^{2}
\;\le\;
\frac{2\,\bigl(\Phi(x_{0})-\Phi^{\inf}\bigr)}{\eta\bigl(1-\tfrac{\eta L}{2}\bigr)K}.
\tag{3}
\]
Consequently $\|G_{\eta}(x_{k})\|\to0$ as $k\to\infty$, and every limit point of $\{x_{k}\}$ is a \emph{stationary point} of $\Phi$, i.e.
\[
0\in \grad f(\bar x)+\partial g(\bar x).
\]
\end{theorem}

\section*{Theoretical Properties}
\begin{enumerate}[label=\arabic*.]
\item \textbf{Descent.} From (3) we obtain the one‑step inequality
\[
\Phi(x_{k})-\Phi(x_{k+1})\;\ge\;\Bigl(\frac{1}{\eta}-\frac{L}{2}\Bigr)\|x_{k+1}-x_{k}\|^{2},
\tag{4}
\]
which shows that the method is a \emph{Lyapunov} descent scheme.
\item \textbf{Hyper‑parameter sensitivity.} The constant $c_{\eta}:=\frac{1}{\eta}-\frac{L}{2}>0$ governs the speed of the descent. As $\eta\!\uparrow\!1/L$, $c_{\eta}\!\downarrow\!0$ and the bound (3) deteriorates, reflecting the usual stepsize trade‑off.
\item \textbf{Comparison to vanilla gradient descent.} When $g\equiv0$, $\prox_{\eta g}$ reduces to the identity and $G_{\eta}(x)=\grad f(x)$. Inequality (3) recovers the classical $O(1/K)$ guarantee for the squared gradient norm of non‑convex smooth optimisation.
\item \textbf{Stability.} The non‑expansiveness (2) of $\prox$ guarantees that the proximal step cannot increase the distance to any reference point, which is essential for the telescoping argument leading to (3).
\end{enumerate}

\section*{Discussion}
The theorem delivers a sublinear rate $O(1/K)$ for the squared norm of the gradient mapping, which is the standard optimal rate for first‑order methods on smooth non‑convex problems. The proof hinges only on the Lipschitz continuity of $\grad f$ and the non‑expansiveness of the proximal operator; no convexity of $f$ is required. Hence the result applies to a broad class of deep‑learning loss functions combined with convex regularisers (e.g. $\ell_{1}$, indicator of a convex set).

\newpage
\section*{Preliminaries (Definitions and Lemmas)}
\begin{lemma}[Descent Lemma]\label{lem:descent}
If $\grad f$ is $L$‑Lipschitz, then for any $x,y\in\R^{d}$
\[
f(y)\le f(x)+\langle\grad f(x),y-x\rangle+\frac{L}{2}\|y-x\|^{2}.
\tag{5}
\]
\end{lemma}
\begin{proof}
Standard integration of the Lipschitz condition; omitted for brevity. \qedhere
\end{proof}
\begin{lemma}[Non‑expansiveness of $\prox$]\label{lem:nonexp}
For a convex $g$ and any $\eta>0$, the proximal operator satisfies (2). \qedhere
\end{lemma}
\begin{lemma}[Relation between gradient mapping and proximal step]\label{lem:gradmap}
For any $x$,
\[
x-\eta G_{\eta}(x)=\prox_{\eta g}\!\bigl(x-\eta\grad f(x)\bigr).
\tag{6}
\]
\end{lemma}
\begin{proof}
Directly from the definition (1). \qedhere
\end{proof}

\section*{Main Proof (Step‑by‑Step Derivation)}
We prove Theorem~\ref{thm:convergence}.  
All steps are numbered for easy reference.

\bigskip
\noindent\textbf{[Step 1]} \emph{Apply the descent lemma to $f$ at $x_{k}$ and $x_{k+1}$.}  
Using \eqref{eq:update} and Lemma~\ref{lem:gradmap},
\[
x_{k+1}=x_{k}-\eta G_{\eta}(x_{k}).
\]
Hence, with $y=x_{k+1}$ in (5),
\[
f(x_{k+1})\le f(x_{k})+\langle\grad f(x_{k}),x_{k+1}-x_{k}\rangle+\frac{L}{2}\|x_{k+1}-x_{k}\|^{2}.
\tag{7}
\]

\noindent\textbf{[Step 2]} \emph{Rewrite the inner product using $x_{k+1}=x_{k}-\eta G_{\eta}(x_{k})$.}
\[
\langle\grad f(x_{k}),x_{k+1}-x_{k}\rangle
= -\eta\langle\grad f(x_{k}),G_{\eta}(x_{k})\rangle .
\tag{8}
\]

\noindent\textbf{[Step 3]} \emph{Use the definition of $G_{\eta}$ to obtain a three‑term identity.}  
From (6),
\[
G_{\eta}(x_{k})=\frac{1}{\eta}\bigl(x_{k}-x_{k+1}\bigr).
\]
Thus,
\[
\langle\grad f(x_{k}),G_{\eta}(x_{k})\rangle
= \frac{1}{\eta}\langle\grad f(x_{k}),x_{k}-x_{k+1}\rangle .
\tag{9}
\]

\noindent\textbf{[Step 4]} \emph{Combine (7)–(9).}
\[
f(x_{k+1})\le f(x_{k})-\langle\grad f(x_{k}),x_{k}-x_{k+1}\rangle
+\frac{L}{2}\|x_{k+1}-x_{k}\|^{2}.
\tag{10}
\]

\noindent\textbf{[Step 5]} \emph{Add the regulariser $g$ and use the optimality condition of the proximal step.}  
By definition of the proximal operator,
\[
x_{k+1}= \prox_{\eta g}(x_{k}-\eta\grad f(x_{k}))
\quad\Longleftrightarrow\quad
0\in \partial g(x_{k+1})+\frac{1}{\eta}(x_{k+1}-x_{k}+\eta\grad f(x_{k})).
\]
Rearranging,
\[
-\grad f(x_{k})\in \partial g(x_{k+1})+\frac{1}{\eta}(x_{k+1}-x_{k}).
\tag{11}
\]
Taking the inner product of (11) with $x_{k+1}-x_{k}$ yields
\[
-\langle\grad f(x_{k}),x_{k+1}-x_{k}\rangle
\ge g(x_{k})-g(x_{k+1})+\frac{1}{\eta}\|x_{k+1}-x_{k}\|^{2}.
\tag{12}
\]
(The inequality follows from the subgradient inequality for convex $g$:
$g(x_{k})\ge g(x_{k+1})+\langle s, x_{k}-x_{k+1}\rangle$ for any $s\in\partial g(x_{k+1})$.)

\noindent\textbf{[Step 6]} \emph{Insert (12) into (10).}
\[
\begin{aligned}
f(x_{k+1}) &\le f(x_{k})+\Bigl[g(x_{k})-g(x_{k+1})+\frac{1}{\eta}\|x_{k+1}-x_{k}\|^{2}\Bigr]
+\frac{L}{2}\|x_{k+1}-x_{k}\|^{2} \\
&= \Phi(x_{k})-\Phi(x_{k+1})+\Bigl(\frac{1}{\eta}+\frac{L}{2}\Bigr)\|x_{k+1}-x_{k}\|^{2}.
\end{aligned}
\tag{13}
\]

\noindent\textbf{[Step 7]} \emph{Rearrange (13) to obtain the fundamental descent inequality.}
\[
\Phi(x_{k})-\Phi(x_{k+1})
\;\ge\;\Bigl(\frac{1}{\eta}-\frac{L}{2}\Bigr)\|x_{k+1}-x_{k}\|^{2}.
\tag{14}
\]
Because $\eta<1/L$, the coefficient $c_{\eta}:=\frac{1}{\eta}-\frac{L}{2}>0$.

\noindent\textbf{[Step 8]} \emph{Express $\|x_{k+1}-x_{k}\|$ through the gradient mapping.}  
From (6),
\[
\|x_{k+1}-x_{k}\|=\eta\|G_{\eta}(x_{k})\|.
\tag{15}
\]

\noindent\textbf{[Step 9]} \emph{Plug (15) into (14).}
\[
\Phi(x_{k})-\Phi(x_{k+1})
\;\ge\;c_{\eta}\,\eta^{2}\|G_{\eta}(x_{k})\|^{2}
\;=\;\Bigl(\frac{1}{\eta}-\frac{L}{2}\Bigr)\eta^{2}\|G_{\eta}(x_{k})\|^{2}.
\tag{16}
\]

\noindent\textbf{[Step 10]} \emph{Sum (16) from $k=0$ to $K-1$ and telescope.}
\[
\sum_{k=0}^{K-1}\bigl[\Phi(x_{k})-\Phi(x_{k+1})\bigr]
= \Phi(x_{0})-\Phi(x_{K})
\;\ge\;c_{\eta}\eta^{2}\sum_{k=0}^{K-1}\|G_{\eta}(x_{k})\|^{2}.
\tag{17}
\]

\noindent\textbf{[Step 11]} \emph{Use the lower bound $\Phi(x_{K})\ge\Phi^{\inf}$.}
\[
\Phi(x_{0})-\Phi^{\inf}\;\ge\;c_{\eta}\eta^{2}\sum_{k=0}^{K-1}\|G_{\eta}(x_{k})\|^{2}.
\tag{18}
\]

\noindent\textbf{[Step 12]} \emph{Derive a bound on the minimal gradient‑mapping norm.}
Since the average of non‑negative numbers dominates the minimum,
\[
\min_{0\le k\le K-1}\|G_{\eta}(x_{k})\|^{2}
\;\le\;\frac{1}{K}\sum_{k=0}^{K-1}\|G_{\eta}(x_{k})\|^{2}
\;\le\;\frac{\Phi(x_{0})-\Phi^{\inf}}{c_{\eta}\eta^{2}K}.
\tag{19}
\]

\noindent\textbf{[Step 13]} \emph{Insert $c_{\eta}= \frac{1}{\eta}-\frac{L}{2}$ to obtain (3).}
\[
\boxed{
\min_{0\le k\le K-1}\|G_{\eta}(x_{k})\|^{2}
\;\le\;
\frac{2\bigl(\Phi(x_{0})-\Phi^{\inf}\bigr)}
{\eta\bigl(1-\tfrac{\eta L}{2}\bigr)K}
}
\]
which completes the proof. \qed

\section*{Rate Derivation (With Constants)}
From (3) we have the explicit $O(1/K)$ decay:
\[
\|G_{\eta}(x_{k})\| = O\!\Bigl(\frac{1}{\sqrt{k}}\Bigr).
\]
If the composite objective satisfies the Polyak–Łojasiewicz (PL) inequality,
\[
\frac{1}{2}\|G_{\eta}(x)\|^{2}\ge\mu\bigl(\Phi(x)-\Phi^{\inf}\bigr),\qquad \mu>0,
\]
then (14) together with the PL condition yields a linear rate
\[
\Phi(x_{k})-\Phi^{\inf}\le (1-\eta\mu)^{k}\bigl(\Phi(x_{0})-\Phi^{\inf}\bigr).
\]

\section*{Special Cases}
\begin{itemize}
\item \textbf{Pure gradient descent ($g\equiv0$).} $G_{\eta}(x)=\grad f(x)$ and (3) reduces to the standard non‑convex gradient bound.
\item \textbf{Indicator of a convex set $\mathcal C$.} $g=\iota_{\mathcal C}$, $\prox_{\eta g}$ becomes the Euclidean projection onto $\mathcal C$. The theorem then guarantees that projected gradient descent converges to a stationary point of the constrained problem.
\item \textbf{Lasso regularisation.} $g(x)=\lambda\|x\|_{1}$, $\prox_{\eta g}$ is the soft‑thresholding operator. The same $O(1/K)$ guarantee holds for the sparse optimisation problem.
\end{itemize}

\end{document}